\section*{الفصل \ref{ch:g11:minerals} --- المعادن}

\begin{solution}{exo:g11:minerals:1}
\omterm{def:g11:minerals:cleavage}{الانفصام} هو نزوع المعدن إلى الانكسار على امتداد مستويات مسطحة، موازية
لاتجاهات تكون فيها الروابط بين ذراته أضعف؛ \omterm{def:g11:minerals:cleavage}{والمكسر} هو
الانكسار على امتداد سطوح غير منتظمة، في معدن لا مستوى ضعف فيه.
ولا \omterm{def:g11:minerals:cleavage}{انفصام} للكوارتز: فهو يُظهر \omterm{def:g11:minerals:cleavage}{مكسرًا} محاريًا، ذا سطوح منحنية ملساء
كالزجاج المكسور. وللميكا \omterm{def:g11:minerals:cleavage}{انفصام} واحد تام: فهي تتقشر
صفائح.
\end{solution}

\begin{solution}{exo:g11:minerals:2}
أثر من عنصر آخر يكفي لتغيير لون الكوارتز، الذي
قد يكون شفافًا أو أبيض أو أسود أو ورديًا أو أزرق أو بنفسجيًا؛ أما الأوليفين فأخضر في الغالب
دائمًا. والأجدر بالثقة: \omterm{def:g11:minerals:cleavage}{الانفصام} (أو \omterm{def:g11:minerals:cleavage}{المكسر})، والصلادة، \omterm{def:g11:minerals:lustre}{والبريق}،
\omterm{def:g11:minerals:habit}{والهيئة البلورية} للبلورات الحرة، والكثافة.
\end{solution}

\begin{solution}{exo:g11:minerals:3}
(1) الأوليفين. (2) البيوتيت. (3) الأورثوكلاز (فلسبار: \omterm{def:g11:minerals:cleavage}{انفصامان}، وصلادة
6). (4) الكوارتز (لا \omterm{def:g11:minerals:cleavage}{انفصام}، وصلادة 7).
\end{solution}

\begin{solution}{exo:g11:minerals:4}
(1) خمس ذرات: سيليكون واحدة في المركز وأربع ذرات أكسجين عند رؤوس
رباعي أوجه. (2) $+4 + 4 \times (-2) = -4$: فالمجموعة هي \ce{SiO4^4-}.
\end{solution}

\begin{solution}{exo:g11:minerals:5}
الأوليفين: رباعيات أوجه منفردة. البيروكسينات: سلاسل مفردة. الأمفيبولات: سلاسل
مزدوجة. الميكا: صفائح. الفلسبارات والكوارتز: هياكل.
\end{solution}

\begin{solution}{exo:g11:minerals:6}
(1) $M(\ce{SiO2}) = 28.085 + 2 \times 15.999 = \qty{60.083}{g/mol}$؛
والأكسجين: $31.998/60.083 \approx \qty{53.3}{\percent}$.
(2) $M(\ce{Mg2SiO4}) = 2 \times 24.305 + 28.085 + 4 \times 15.999 =
\qty{140.691}{g/mol}$؛ والأكسجين: $63.996/140.691 \approx \qty{45.5}{\percent}$.
(3) القشرة، بنسبة \qty{46}{\percent}، من الرتبة نفسها: فهي مصنوعة كلها تقريبًا
من معادن غنية بالأكسجين، هي \omterm{def:g11:minerals:silicate}{السيليكات}.
\end{solution}

\begin{solution}{exo:g11:minerals:7}
للحبيبة البيضاء \omterm{def:g11:minerals:cleavage}{انفصامان} وصلادة بين 5.5 و6.5: فهي
فلسبار (6). والحبيبة الرمادية أصلد من 6.5 وتُظهر \omterm{def:g11:minerals:cleavage}{مكسرًا} منحنيًا
محاريًا بلا \omterm{def:g11:minerals:cleavage}{انفصام}: فهي كوارتز (7).
\end{solution}

\begin{solution}{exo:g11:minerals:8}
(1) $4 - 2/2 = 3$ ذرات أكسجين لكل ذرة سيليكون: \ce{SiO3}. (2) يضم \ce{(Si,Al)2O6}
ذرتي سيليكون (أو ألمنيوم) لست ذرات أكسجين، أي ثلاثًا لكل ذرة سيليكون: فهو
\ce{SiO3} مرتين.
\end{solution}

\begin{solution}{exo:g11:minerals:9}
$\rho = 0.55 \times 2.69 + 0.35 \times 3.4 + 0.10 \times 3.32 = 1.480 + 1.190
+ 0.332 \approx \qty{3.00}{g/cm^3}$. وحجم المكعب
\qty{1000}{cm^3}، فكتلته نحو \qty{3.00}{kg}.
\end{solution}

\begin{solution}{exo:g11:minerals:10}
(1) $0.75 \times 3.32 + 0.25 \times 3.4 = 2.49 + 0.85 = \qty{3.34}{g/cm^3}$.
(2) $(3.34 - 3.00)/3.00 \approx 0.11$: أي أكثف بنحو \qty{11}{\percent}.
(3) بازلت \omterm{def:g11:ocean-ridges:oceanic-lithosphere}{القشرة المحيطية} وغابروها أخف من
\omterm{def:g11:ocean-ridges:peridotite}{بيريدوتيت} الوشاح: فيبقيان فوقه.
\end{solution}

\begin{solution}{exo:g11:minerals:11}
(1) البلاجيوكليز $168/400 = 0.42$، والبيروكسين $0.33$، والأوليفين $0.15$، والثقوب
$0.10$. (2) القضية تصح لصخر بلا مسام؛ وهنا عُشر
الحجم فارغ. (3) في الجزء الصلب تكون الكسور $0.42/0.90
\approx 0.467$ و$0.367$ و$0.167$:
$\rho_{\text{صلب}} \approx 0.467 \times 2.69 + 0.367 \times 3.4 + 0.167
\times 3.32 \approx \qty{3.06}{g/cm^3}$. ومع ثقوبه، التي تضيف حجمًا
بلا كتلة، تكون كثافة الصخر $0.42 \times 2.69 + 0.33 \times 3.4 +
0.15 \times 3.32 \approx \qty{2.75}{g/cm^3}$، أي $0.9 \times 3.06$.
\end{solution}

\begin{solution}{exo:g11:minerals:12}
(1) الأكسجين $464/16.0 = 29.0$~\unit{mol}؛ والسيليكون $283/28.1 \approx 10.1$؛ والألمنيوم
$84/27.0 \approx 3.1$؛ والحديد $52/55.8 \approx 0.93$؛ والكالسيوم $46/40.1 \approx
1.15$؛ والمغنيسيوم $28/24.3 \approx 1.15$؛ والصوديوم $23/23.0 = 1.0$؛ والبوتاسيوم
$15/39.1 \approx 0.38$. (2) المجموع نحو $46.8$~\unit{mol} من الذرات، منها
$29.0/46.8 \approx 0.62$: فذرتان من كل ثلاث ذرات تقريبًا ذرتا أكسجين.
(3) أيونات الأكسجين الكبيرة تشكّل معظم الذرات، وتشغل حصة أكبر حتى من
الحجم: فالقشرة، على مقياس الذرات، رصّ من أيونات الأكسجين، تقع
أيونات السيليكون والفلزات الصغيرة في الفراغات بينها.
\end{solution}

\begin{solution}{exo:g11:minerals:13}
(1) متوسط عدد الرؤوس المشتركة $(2 + 3)/2 = 2.5$، فيكون عدد ذرات الأكسجين
$4 - 2.5/2 = 2.75$ لكل ذرة سيليكون. (2) أربع ذرات سيليكون لإحدى عشرة ذرة أكسجين:
\ce{Si4O11}. (3) يضم \ce{Si7AlO22} ثمانية رباعيات أوجه (سبعة للسيليكون، وواحد
للألمنيوم) مقابل 22 ذرة أكسجين، أي $22/8 = 2.75$ لكل رباعي أوجه، أو
\ce{Si4O11} مرتين. ومجموعتا \ce{OH} ليستا جزءًا من رباعيات الأوجه.
\end{solution}

\begin{solution}{exo:g11:minerals:14}
$0.20 \times 2.62 + 0.10 \times 2.56 + 0.50 \times 2.65 + 0.10 \times 3.09 +
0.10 \times 3.23 = 0.524 + 0.256 + 1.325 + 0.309 + 0.323 \approx
\qty{2.74}{g/cm^3}$، أي فوق \qty{2.7}{g/cm^3}؛ أما غرانيت المثال،
الأغنى بالكوارتز وفلسبار البوتاسيوم، فنحو \qty{2.59}{g/cm^3}، أي تحت
\qty{2.6}{g/cm^3}. فصخور الغرب أكثف لأنها تحوي مزيدًا من
البلاجيوكليز ومزيدًا من البيوتيت والهورنبلند، الغنيين بالحديد والمغنيسيوم؛ أما
الغرانيتات الغنية بالسيليكا فتحوي مزيدًا من الكوارتز وفلسبار البوتاسيوم، وهما أخف
المعادن المكوّنة للصخور.
\end{solution}

\begin{solution}{exo:g11:minerals:15}
(1) القشرة العليا: $66.6 \times 0.467 \approx \qty{31.1}{\percent}$ من السيليكون؛
والقشرة كلها: $60.6 \times 0.467 \approx \qty{28.3}{\percent}$. (2) ينبغي أن تحتوي القشرة
العليا، الأغنى بالسيليكا، على مزيد من الكوارتز والفلسبارات، وعلى قدر أقل من
المعادن الغنية بالحديد والمغنيسيوم؛ والقشرة الأعمق على العكس. (3) الكوارتز
والفلسبارات أخف المعادن المكوّنة للصخور: فينبغي أن تكون القشرة العليا
أخف من القشرة الأعمق.
\end{solution}

\begin{solution}{pb:g11:minerals:1}
\textbf{1.} الكوارتز: لا \omterm{def:g11:minerals:cleavage}{انفصام}، \omterm{def:g11:minerals:cleavage}{ومكسر} محاري، \omterm{def:g11:minerals:lustre}{وبريق} زجاجي، وهو أصلد
من المسمار (7).

\textbf{2.} الأورثوكلاز، وهو فلسبار بوتاسيوم: \omterm{def:g11:minerals:cleavage}{انفصامان}، وصلادة 6 على
سلم موس.

\textbf{3.} البلاجيوكليز (هنا الأوليغوكلاز)، وهو الفلسبار الذي يحتوي على الصوديوم
والكالسيوم.

\textbf{4.} D هي البيوتيت، وE المسكوفيت: نوعان من الميكا. ورباعيات أوجههما تكوّن
صفائح، والروابط بين الصفائح ضعيفة: فتنفصل على امتداد هذا \omterm{def:g11:minerals:cleavage}{الانفصام}
التام الوحيد.

\textbf{5.} الكوارتز والفلسباران (A وB وC) هياكل؛ والبيوتيت
(D) يحتوي على الحديد والمغنيسيوم.

\textbf{6.} الكوارتز $150/500 = 0.30$؛ والأورثوكلاز $0.35$؛ والأوليغوكلاز $0.25$؛
والبيوتيت $0.06$؛ والمسكوفيت $0.04$.

\textbf{7.} حين تكون الحبيبات مبعثرة بلا نظام، فإن نصيب النقاط الواقعة
على معدن يقدّر نصيب المساحة التي يغطيها من الوجه، وهذا يساوي
نصيب الحجم الذي يملؤه.

\textbf{8.} $0.35 + 0.25 = 0.60$: أي \qty{60}{\percent} فلسبار.

\textbf{9.} بمقدار $1/500 = 0.002$، أي \qty{0.2}{\percent}.

\textbf{10.} نعم: فلسبار وكوارتز وميكا.

\textbf{11.} $\rho = \phi_1 \rho_1 + \phi_2 \rho_2 + \dots + \phi_k \rho_k$،
مع $\phi_1 + \dots + \phi_k = 1$.

\textbf{12.} الكوارتز $0.30 \times 2.62 = 0.786$؛ والأورثوكلاز $0.35 \times 2.56 =
0.896$؛ والأوليغوكلاز $0.25 \times 2.65 = 0.6625$؛ والبيوتيت $0.06 \times 3.09
\approx 0.185$؛ والمسكوفيت $0.04 \times 2.82 \approx 0.113$ (كلها
بوحدة \unit{g/cm^3}).

\textbf{13.} $0.786 + 0.896 + 0.6625 + 0.185 + 0.113 \approx
\qty{2.64}{g/cm^3}$.

\textbf{14.} الأورثوكلاز، بمقدار $0.896$، لأنه الأوفر. أما البيوتيت
والمسكوفيت، بكثافتي \qty{3.09}{} و\qty{2.82}{g/cm^3}، فهما أكثف من
المعادن الأخرى ويرفعان كثافة الصخر فوق كثافة الكوارتز
والفلسبار وحدهما.

\textbf{15.} \qty{2640}{kg/m^3}: فكتلة \qty{1}{m^3} من الغرانيت
نحو \qty{2.64}{t}.

\textbf{16.} $100 \times 2.643 \approx \qty{264.3}{g}$.

\textbf{17.} تزيح الكتلة الصخرية \qty{100}{cm^3} من الماء، كتلتها
\qty{100}{g}: فيقرأ الميزان $264.3 - 100 = \qty{164.3}{g}$.

\textbf{18.} $264.3/(264.3 - 164.3) = 264.3/100 \approx \qty{2.64}{g/cm^3}$.

\textbf{19.} بين الاثنين: أخف من المتوسط البالغ \qty{2.69}{g/cm^3}،
وفوق \qty{2.6}{g/cm^3} التي للغرانيتات الأغنى بالسيليكا والبوتاسيوم.
ومزيد من البلاجيوكليز ومزيد من البيوتيت أو الهورنبلند، وقدر أقل من الكوارتز
وفلسبار البوتاسيوم، يجعل الغرانيت أكثف.

\textbf{20.} كثافة الغرانيت نحو
\textbf{\qty{2.64}{g/cm^3}}.
\end{solution}
